\chapter{Dysmenorrhoea}
\index{Dysmenorrhoea}

\section*{Definition and Overview}
Dysmenorrhoea is the medical term for painful periods, describing cramping pain in the lower abdomen that occurs around menstruation. It is one of the most common gynaecological problems and a major cause of absenteeism from school and work. Primary dysmenorrhoea is period pain with no underlying pelvic disease, typically beginning in adolescence, while secondary dysmenorrhoea is pain caused by a condition such as endometriosis, adenomyosis, or fibroids. Severe period pain is treatable, and women should not be expected to simply cope with it.

\section*{Epidemiology}
Painful periods affect a large proportion of menstruating women worldwide, with severe pain reported by a significant minority. The condition is most prevalent among adolescents and young women and commonly improves with age. Secondary dysmenorrhoea becomes more common in the 30s and 40s. Many women manage without medical help, and a substantial number do not know that effective treatment exists.

\section*{Aetiology and Pathophysiology}
Period pain results from uterine muscle contractions that are driven by chemical messengers called prostaglandins.
\begin{itemize}
    \item \textbf{Prostaglandins:} levels rise in the uterine lining in the days before menstruation and cause strong, painful uterine contractions and reduced blood flow to the uterus
    \item \textbf{Primary dysmenorrhoea:} pain without an underlying cause, usually starting one to three years after the first period
    \item \textbf{Secondary dysmenorrhoea:} pain related to conditions including endometriosis, adenomyosis, fibroids, and pelvic infection
    \item \textbf{Risk factors:} early menarche, heavy or long periods, smoking, family history, and stress
    \item \textbf{Chronic pain mechanisms:} in some women, pelvic pain pathways become sensitised and pain continues beyond menstruation
\end{itemize}

\section*{Clinical Features}
The hallmark is cramping pain in the lower abdomen, but the full picture varies.
\begin{itemize}
    \item Cramping, colicky pain in the lower abdomen, often radiating to the lower back and inner thighs
    \item Pain beginning just before or with the onset of bleeding and lasting one to three days
    \item Nausea, vomiting, diarrhoea, and headache
    \item Fatigue, dizziness, and difficulty concentrating
    \item Pain that eases as the period progresses
    \item In secondary dysmenorrhoea, pain may start earlier in the cycle, last longer, and be accompanied by heavy bleeding or deep pain during intercourse
\end{itemize}

\section*{Differential Diagnosis}
Other causes of pelvic pain must be considered, particularly in secondary dysmenorrhoea.
\begin{itemize}
    \item Endometriosis, with pain through the cycle, painful intercourse, and infertility
    \item Adenomyosis, with heavy, painful periods, often in women in their 30s and 40s
    \item Uterine fibroids and polyps
    \item Pelvic inflammatory disease
    \item Ovarian cysts and ovarian pathology
    \item Gastrointestinal conditions such as irritable bowel syndrome
    \item Urinary tract problems
    \item Ectopic pregnancy or miscarriage in women of reproductive age
\end{itemize}

\section*{When to Seek Medical Care}
Women should see a GP if period pain is severe, not controlled by over-the-counter medicines, or affecting daily life, or if it is accompanied by heavy bleeding, irregular bleeding, pain between periods, or pain during sex.

\textbf{Contact your local emergency services for:}
\begin{itemize}
    \item Sudden severe pelvic pain, especially with the possibility of pregnancy
    \item Heavy bleeding with faintness, dizziness, or collapse
    \item Fever with pelvic pain
    \item Shoulder tip pain with pelvic pain and a positive pregnancy test, suggesting ectopic pregnancy
\end{itemize}

\section*{Diagnosis and Investigations}
The aim is to distinguish primary from secondary dysmenorrhoea.
\begin{itemize}
    \item \textbf{History:} the timing, nature, and severity of pain and its relationship to the cycle
    \item \textbf{Pelvic examination:} to identify tenderness, masses, or infection
    \item \textbf{Pregnancy test:} to exclude pregnancy
    \item \textbf{Pelvic ultrasound:} to detect fibroids, cysts, and adenomyosis
    \item \textbf{Laparoscopy:} for suspected endometriosis
    \item \textbf{Response to treatment:} improvement with NSAIDs and the contraceptive pill supports primary dysmenorrhoea
\end{itemize}

\section*{Management}
Treatment is effective and should be tailored to the woman's needs and preferences.
\begin{itemize}
    \item \textbf{Anti-inflammatory medicines:} ibuprofen or naproxen taken at the first sign of pain or bleeding
    \item \textbf{Hormonal contraception:} the combined pill, contraceptive implant, or hormonal IUD reduce pain in most women
    \item \textbf{Heat therapy:} a heat pack on the lower abdomen
    \item \textbf{Regular exercise:} physical activity reduces the severity of period pain
    \item \textbf{Relaxation and complementary therapies:} stress reduction, yoga, and acupuncture help some women
    \item \textbf{Treatment of the underlying cause:} for secondary dysmenorrhoea, management of endometriosis, fibroids, or infection
    \item \textbf{Surgical options:} for example, removal of fibroids or treatment of endometriosis
\end{itemize}

\section*{Complications}
\begin{itemize}
    \item Absenteeism from school and work and lost productivity
    \item Reduced quality of life and wellbeing
    \item Delay in diagnosis of underlying conditions such as endometriosis
    \item Chronic pelvic pain in some women
    \item Psychological distress and anxiety around menstruation
\end{itemize}

\section*{Prognosis and Outcomes}
Most women with primary dysmenorrhoea respond well to simple measures, and the pain often lessens with age and after childbirth. Women with secondary dysmenorrhoea generally improve with treatment of the underlying condition. Seeking help early and finding an effective individual treatment are the keys to a good outcome.

\section*{Prevention and Self-Care}
\begin{itemize}
    \item Keep a diary of symptoms and treatments that help
    \item Start anti-inflammatory medicines early, at the first sign of pain
    \item Use heat packs for prompt relief
    \item Stay active with regular exercise
    \item Manage stress with relaxation techniques
    \item Consider hormonal contraception with medical advice
    \item See a doctor if simple measures do not control the pain
\end{itemize}

\section*{Sources and Further Reading}
The following reputable sources provide further information for healthcare students and patients seeking reliable, current advice about painful periods.
\begin{itemize}
    \item Jean Hailes for Women's Health. \textit{Painful periods.} https://www.jeanhailes.org.au/
    \item Healthdirect Australia. \textit{Period pain (dysmenorrhoea).} https://www.healthdirect.gov.au/period-pain
    \item NHS. \textit{Period pain: symptoms and treatment.} https://www.nhs.uk/conditions/period-pain/
    \item Mayo Clinic. \textit{Menstrual cramps.} https://www.mayoclinic.org/diseases-conditions/menstrual-cramps/
    \item Royal Women's Hospital. \textit{Period pain information.} https://www.thewomens.org.au/
    \item MSD Manual. \textit{Dysmenorrhea.} https://www.msdmanuals.com/
\end{itemize}
